\documentclass[11pt,letterpaper]{article}
\usepackage[T1]{fontenc}
\usepackage{lmodern,microtype}
\usepackage[margin=.78in,headheight=15pt,headsep=15pt]{geometry}
\usepackage{mathtools,amssymb,booktabs,array,xcolor}
\usepackage[most]{tcolorbox}
\usepackage{enumitem,fancyhdr,needspace,hyperref}
\definecolor{ink}{HTML}{17243A}
\definecolor{indigo}{HTML}{293D70}
\definecolor{mint}{HTML}{2A8C82}
\definecolor{muted}{HTML}{647084}
\definecolor{mist}{HTML}{EAF5F2}
\color{ink}
\hypersetup{colorlinks=true,linkcolor=mint,urlcolor=mint,
 pdfauthor={Igor Rivin},pdftitle={Probability Problem Set 4: Distributions, Tails, and Coupon Collection}}
\setlength{\parindent}{0pt}
\setlength{\parskip}{4pt}
\clubpenalty=10000
\widowpenalty=10000
\setlist[enumerate]{label=(\alph*),leftmargin=1.8em,itemsep=3pt,topsep=4pt}
\setlength{\abovedisplayskip}{7pt}
\setlength{\belowdisplayskip}{7pt}
\renewcommand{\P}{\mathbb P}
\newcommand{\E}{\mathbb E}
\DeclareMathOperator{\Var}{Var}
\DeclareMathOperator{\Geom}{Geom}
\DeclareMathOperator{\Bin}{Bin}
\DeclareMathOperator{\Pois}{Pois}
\newcommand{\pts}[1]{\hfill\textcolor{muted}{\small[#1 points]}}
\newcommand{\partpts}[1]{\unskip\nobreak\hspace{.3em}\mbox{\textcolor{muted}{\small(#1)}}}
\newcommand{\edition}{\ifsolutions Instructor solutions\else Student handout\fi}
\newcounter{problem}
\newcommand{\problem}[2]{\par\addvspace{11pt}\Needspace{7\baselineskip}\stepcounter{problem}{\bfseries\color{indigo}\theproblem.\ #1}\pts{#2}\par\nobreak}
\newcommand{\topic}[1]{\par\addvspace{9pt}\Needspace{10\baselineskip}{\large\bfseries\color{mint}#1}\par\nobreak}
\newcommand{\studentbreak}{\ifsolutions\else\newpage\fi}
\pagestyle{fancy}
\fancyhf{}
\fancyhead[L]{\small\color{muted}Probability \textbar\ Problem Set 4}
\fancyhead[R]{\small\color{muted}Chance, Data, and Decisions}
\fancyfoot[L]{\footnotesize\color{muted}\edition}
\fancyfoot[R]{\small\color{mint}\thepage}
\renewcommand{\headrulewidth}{.3pt}
\NewDocumentEnvironment{solution}{+b}{\ifsolutions
 \begin{tcolorbox}[enhanced,breakable,colback=mist!45,colframe=mint!45,
  boxrule=.5pt,arc=1mm,left=3mm,right=3mm,top=2mm,bottom=2mm,
  before skip=7pt,after skip=7pt,title={Solution},
  title after break={Solution (continued)},colbacktitle=mist,coltitle=mint,
  fonttitle=\small\bfseries]\small #1\end{tcolorbox}\fi}{}

\begin{document}
{\LARGE\bfseries\color{indigo}Problem Set 4}\hfill{\small\color{mint}\edition}\\[4pt]
{\large Distributions, tails, and coupon collection}\\[4pt]
{\small\color{muted}100 points + 10 optional challenge points \quad $\bullet$ \quad About 4 hours}\par
\ifsolutions
\textbf{Teaching note.} Problems 3--4 deepen the coupon collector; Problems
6--7 develop tail areas and use both sides of a signed payoff.
The challenge introduces unequal coupon probabilities.
Award the indicated subpart points for correct reasoning, allowing follow-through
credit after arithmetic errors. The tail proofs and challenge are guided enrichment.
\else
\textbf{Name:}\ \rule{.46\linewidth}{.4pt}\hfill
\textbf{Due:}\ \rule{.20\linewidth}{.4pt}
\fi
\par\medskip
\textbf{Instructions.} Show your reasoning and state each distribution's
support and parameters. Geometric variables count trials including the
first success; exponential parameters are rates. Write $H_r=\sum_{j=1}^r1/j$
and $H_0=0$. Exact expressions suffice unless decimals are requested.
A calculator or short program is allowed for numerical comparisons.
Independent uniform choices of coupon types are assumed except in the challenge.

\topic{I. Discrete mechanisms}
\problem{A count is not a wait}{12}
A box contains 20 distinguishable gadgets, of which 6 are defective.
\begin{enumerate}
\item Choose one gadget uniformly. Let $I$ indicate a defect. Give its
law, mean, and variance. \partpts{2}
\item Make four independent uniform draws \emph{with replacement}.
For the number $B$ of defective gadgets drawn, give the law, $\P(B=2)$,
mean, and variance. \partpts{3}
\item Choose four gadgets \emph{without replacement}. For the defective
count $D$, give its support, PMF, mean, and variance. Explain the variance
comparison with (b). \partpts{4}
\item Draw independently with replacement until the first defect.
For the draw count $G$, give its PMF, $\P(G>4)$, mean, and variance.
Why is it not binomial? \partpts{3}
\end{enumerate}
\begin{solution}
\begin{enumerate}
\item $I\sim\operatorname{Bernoulli}(3/10)$ on $\{0,1\}$:
$\E[I]=3/10$, $\Var(I)=21/100$.
\item $B\sim\Bin(4,3/10)$ on $\{0,\ldots,4\}$.
$\P(B=2)=\binom42(3/10)^2(7/10)^2=1323/5000$;
$\E[B]=6/5$, $\Var(B)=21/25$.
\item $D$ is hypergeometric with $N=20,K=6,m=4$, supported on
$\{0,\ldots,4\}$, with
$\P(D=k)=\binom6k\binom{14}{4-k}/\binom{20}4$.
$\E[D]=6/5$, $\Var(D)=(21/25)(16/19)=336/475$.
Removing a defect lowers the chance of a defect next time: the draw
indicators have negative covariance, reducing the variance.
\item $G\sim\Geom(3/10)$ on $\{1,2,\ldots\}$,
$\P(G=k)=(7/10)^{k-1}(3/10)$, $\P(G>4)=(7/10)^4$,
$\E[G]=10/3$, $\Var(G)=70/9$.
The number of trials is random; we record the first-success time,
not the successes in a fixed number of trials.
\end{enumerate}
\end{solution}

\problem{Rare errors and a geometric tail}{10}
Each of 2000 independent transmitted bits is corrupted with probability
$0.001$. Let $N$ count corrupted bits.
\begin{enumerate}
\item Give the exact law of $N$ and a suitable Poisson approximation.
Compare the exact and approximate $\P(N\ge2)$ to four decimals. \partpts{4}
\item For $Q\sim\Pois(\lambda)$, derive $\E[Q(Q-1)]=\lambda^2$ by
shifting the PMF sum, and recover its variance. \partpts{3}
\item In an unlimited stream with the same independent error probability,
let $G$ be the bit index of the first error. Use
$\E[G]=\sum_{m=0}^{\infty}\P(G>m)$ to find its mean, and calculate
$\P(G>s+t\mid G>s)$ for integers $s,t\ge0$. \partpts{3}
\end{enumerate}
\begin{solution}
\begin{enumerate}
\item $N\sim\Bin(2000,.001)$; the approximation is $\Pois(2)$.
The probabilities are
$1-(.999)^{2000}-2000(.001)(.999)^{1999}\approx.5941$ and
$1-3e^{-2}\approx.5940$. The opportunities are numerous, independent,
and rare, with mean count $2000(.001)=2$.
\item For $\lambda>0$,
$\sum_{k\ge2}k(k-1)e^{-\lambda}\lambda^k/k!
=\lambda^2e^{-\lambda}\sum_{j\ge0}\lambda^j/j!=\lambda^2$.
Thus $\E[Q^2]=\lambda^2+\lambda$ and $\Var(Q)=\lambda$.
At $\lambda=0$ the variable is constantly zero.
\item $\P(G>m)=.999^m$, so $\E[G]=1/(1-.999)=1000$ bits.
The conditional probability is $.999^{s+t}/.999^s=.999^t$:
past failures do not change the extra wait.
\end{enumerate}
\end{solution}

\studentbreak
\topic{II. What the coupon-collector mean leaves out}
\problem{The last few stickers}{14}
An album has $n$ equally likely sticker types. Let $T_{n,r}$ be the number
of purchases needed to see $r$ distinct types, $1\le r\le n$, and
let $T_n=T_{n,n}$.
\begin{enumerate}
\item Define stages $W_k$, for $0\le k<n$, identify their laws, and
justify their mutual independence from the conditional distribution after
any completed history. \partpts{4}
\item Derive $\E[T_{n,r}]$ and $\Var(T_{n,r})$ as finite sums. Simplify
the expectation with harmonic numbers. \partpts{4}
\item For $n=100$, compare the expected purchases for the first 90 types
with the expected \emph{additional} purchases for the last 10. Give one
decimal place. Why do the last 10 take so long? \partpts{3}
\item For $n=6$, find $\E[T_6]$ and $\Var(T_6)$ exactly. If exactly two
faces are still missing, what is the expected remaining number of rolls?
Why is $T_6$ not geometric? \partpts{3}
\end{enumerate}
\begin{solution}
\begin{enumerate}
\item $W_k$ counts additional draws, including the next new type, after
$k$ distinct types have appeared. $W_k\sim\Geom((n-k)/n)$.
Given any particular completed history, future draws are independent and
uniform; the chance of a new type stays $(n-k)/n$ until success.
The conditional stage law depends only on $k$, not the previous stage
lengths or labels. Successively multiplying these conditional PMFs gives
the product of the marginal PMFs, proving stage independence. $W_0=1$.
\item $T_{n,r}=\sum_{k=0}^{r-1}W_k$. Hence
\[
\E[T_{n,r}]=n\sum_{j=n-r+1}^n\frac1j=n(H_n-H_{n-r}),\qquad
\Var(T_{n,r})=\sum_{j=n-r+1}^n\left(\frac{n^2}{j^2}-\frac nj\right).
\]
Linearity alone proves the mean; independence is needed for the variance.
\item First 90: $100(H_{100}-H_{10})\approx225.8$.
Last 10: $100H_{10}\approx292.9$. With $j$ types missing, a new one
has probability $j/100$, so its expected wait is $100/j$.
The very last type alone costs 100 purchases on average.
\item $\E[T_6]=6H_6=147/10$;
$\Var(T_6)=36\sum_{j=1}^6j^{-2}-6H_6=3899/100$.
Two missing faces require $6/2+6/1=9$ more rolls on average.
Completion cannot occur before six rolls, and finding a new face changes
the success probability; $T_6$ is a sum of stages, not one geometric wait.
\end{enumerate}
\end{solution}

\problem{A deadline is not an expectation}{14}
Let $M_m$ count types missing after $m\ge1$ purchases, with $n\ge2$ types.
Write $J_i=1$ if type $i$ is missing, $q=(1-1/n)^m$, and
$b=(1-2/n)^m$.
\begin{enumerate}
\item Compute $\E[M_m]$ and $\E[J_iJ_j]$ for $i\ne j$. Deduce
\[
\Var(M_m)=nq(1-q)+n(n-1)(b-q^2).
\]
Are missing-type indicators independent? \partpts{4}
\item Derive the exact completion CDF by inclusion--exclusion:
$\P(T_n\le m)=\sum_{j=0}^n(-1)^j\binom nj(1-j/n)^m$. \partpts{3}
\item For six types, compute $\P(T_6\le15)$ and $\E[M_{15}]$ to four
decimals. Find the smallest integer $m$ with $\P(T_6\le m)\ge.95$;
report probabilities at $m-1$ and $m$. \partpts{4}
\item Prove $\P(T_n>m)\le ne^{-m/n}$. For 100 types choose an integer
$m$ guaranteeing at least $95\%$ completion by this bound. Does the bound
establish the smallest possible deadline? \partpts{3}
\end{enumerate}
\begin{solution}
\begin{enumerate}
\item $M_m=\sum_iJ_i$, so $\E[M_m]=nq$.
Two specified types are both absent with probability $b$, since each
draw avoids them with probability $1-2/n$. Expanding the variance of
a sum gives the stated formula. For $m\ge1$, $b<q^2$ because
$1-2/n<(1-1/n)^2$. Thus the distinct indicators have negative covariance
and are dependent. (At $m=0$ all indicators are constant.)
\item Completion is the complement of the union of missing-type events.
A specified set of $j$ types is absent with probability $(1-j/n)^m$;
there are $\binom nj$ such sets. Inclusion--exclusion gives the result.
At $m=0$ the CDF is zero; handle it directly to avoid $0^0$ conventions.
\item $\P(T_6\le15)\approx.6442$, $\E[M_{15}]=6(5/6)^{15}\approx.3894$.
The smallest deadline is $m=27$, since
$\P(T_6\le26)\approx.9480$ and $\P(T_6\le27)\approx.9566$.
The CDF is nondecreasing. Fifteen rolls are close to the mean 14.7,
yet completion is far from certain.
\item The union bound gives $\P(T_n>m)\le n(1-1/n)^m\le ne^{-m/n}$.
Take $m=\lceil100(\log100+\log20)\rceil=761$.
This is a sufficient deadline; the bound does not determine the exact
minimum. At the chosen $m$, its failure bound is about $.04955$.
\end{enumerate}
\end{solution}

\studentbreak
\topic{III. Densities and tail areas}
\problem{A density that exceeds one}{12}
Suppose $f_X(x)=c x^2$ for $0\le x\le2$, and zero elsewhere.
\begin{enumerate}
\item Normalize the density and write its full piecewise CDF. \partpts{3}
\item Find $\P(1/2<X\le1)$, $\E[X]$, and $\Var(X)$. \partpts{3}
\item Find the median and $\P(X=1)$. Explain why $f_X(2)>1$ is allowed.
\partpts{3}
\item Let $Y=(X/2)^3$. Find $F_Y(y)$ by rewriting $\{Y\le y\}$ as an
event about $X$, and identify the distribution of $Y$. \partpts{3}
\end{enumerate}
\begin{solution}
\begin{enumerate}
\item $1=c\int_0^2x^2\,dx=8c/3$, so $c=3/8$.
$F_X(x)=0$ for $x<0$, $x^3/8$ for $0\le x\le2$, and 1 for $x>2$.
\item $\P(1/2<X\le1)=7/64$;
$\E[X]=\int_0^2(3/8)x^3\,dx=3/2$;
$\E[X^2]=\int_0^2(3/8)x^4\,dx=12/5$.
Thus $\Var(X)=12/5-9/4=3/20$.
\item The median is $2\cdot2^{-1/3}$, solving $F_X(x)=1/2$.
$\P(X=1)=0$. Density is probability per unit length; only its integrated
area must be at most one on an interval. Here $f_X(2)=3/2$.
\item For $0\le y\le1$,
$F_Y(y)=\P(X\le2y^{1/3})=(2y^{1/3})^3/8=y$.
The CDF is zero below 0 and one above 1, so $Y\sim\operatorname{Uniform}(0,1)$.
\end{enumerate}
\end{solution}

\problem{Tail areas and the suspiciously generous prize}{12}
For a real variable $Z$, write $F(z)=\P(Z\le z)$ and assume
$\E[|Z|]<\infty$ in (a)--(b).
\begin{enumerate}
\item For $Z_+=\max(Z,0)$ and $Z_-=\max(-Z,0)$, use
$y=\int_0^\infty\mathbf1_{\{y>t\}}\,dt$ for $y\ge0$ to show
\[
\E[Z]=\int_0^\infty(1-F(t))\,dt-\int_{-\infty}^0F(t)\,dt.
\]
You may interchange expectation and nonnegative integrals. Explain why
using $\P(Z<t)$ instead of $\P(Z\le t)$ does not change an integral.
\partpts{3}
\item What does the \emph{sum} of these two nonnegative tail areas compute?
Check both identities for the constant payoff $Z=-2$. \partpts{3}
\item A carnival advertises a prize $R$ (in dollars) with CDF zero for
$x<1$ and $F_R(x)=1-1/x$ for $x\ge1$. Show that this is a valid
continuous law, find its density, and use a tail integral to find its
mean. Is there a finite fair entry fee for the uncapped game? \partpts{3}
\item The small print caps the payout at $B>1$ dollars: $C=\min(R,B)$.
Use its tail to find the fair entry fee. What probability mass is at
$B$? Evaluate the fee for $B=100$. \partpts{3}
\end{enumerate}
\begin{solution}
\begin{enumerate}
\item Taking expectations in the indicator identity gives
$\E[Z_+]=\int_0^\infty\P(Z>t)\,dt$ and
$\E[Z_-]=\int_0^\infty\P(Z<-t)\,dt$.
A distribution has at most countably many atoms: the jump locations
can change the integrand only at a set of length zero.
Substitute $u=-t$ in the second integral and use
$Z=Z_+-Z_-$. The integrability assumption makes both areas finite.
\item The sum equals $\E[Z_++Z_-]=\E[|Z|]$.
For $Z=-2$, the negative-side area is 2 and the positive-side area is 0;
the signed mean is $-2$, while the sum is 2.
\item The CDF is continuous and nondecreasing, with limits zero and one.
Its density is $x^{-2}$ on $x>1$, zero below 1, and has total area one.
The tail equals 1 on $[0,1)$ and $1/t$ for $t\ge1$. Thus
$\E[R]=1+\int_1^\infty dt/t=+\infty$.
There is no finite fee making expected net payoff zero. This is an
infinite positive mean, not an undefined difference of two infinite areas.
\item $\P(C>t)=1$ for $0\le t<1$, $1/t$ for $1\le t<B$, and zero
for $t\ge B$. Therefore $\E[C]=1+\log B$ dollars, the fair entry fee.
$\P(C=B)=\P(R\ge B)=1/B$: the capped law is a mixture of a density
below $B$ and an atom at $B$. At $B=100$, the fair fee is
$1+\log100\approx\$5.61$, despite the uncapped infinite mean.
\end{enumerate}
\end{solution}

\studentbreak
\topic{IV. Continuous mechanisms and a signed payoff}
\problem{The auction for the last slice}{12}
At a pizza party, $n\ge1$ hungry mathematicians independently submit bids
$U_1,\ldots,U_n$, each uniform on $[0,1]$ dollars. The seller receives
the \emph{highest bid} $M=\max_iU_i$, even if it is small. The slice cost
$c$ dollars, where $0<c<1$, so the seller's net payoff is $Z=M-c$.
\begin{enumerate}
\item Find the full CDF of $M$, and then of $Z$, using independence.
\partpts{3}
\item Compute separately the negative-tail area
$L=\int_{-\infty}^0F_Z(z)\,dz$ and positive-tail area
$G=\int_0^\infty(1-F_Z(z))\,dz$. Use $G-L$ to find the mean profit
\emph{without differentiating the CDF}. \partpts{4}
\item Find the loss probability and $\E[|Z|]$. Explain why the latter
need not equal the absolute value of the mean profit. \partpts{3}
\item For $n=3,c=3/4$, calculate the mean profit, loss probability, and
mean absolute payoff. What does this say about calling a game ``fair''?
\partpts{2}
\end{enumerate}
\begin{solution}
\begin{enumerate}
\item On $0\le x\le1$, $F_M(x)=\prod_i\P(U_i\le x)=x^n$;
it is zero below 0 and one above 1.
$F_Z(z)=0$ below $-c$, $(z+c)^n$ on $[-c,1-c]$, and 1 above $1-c$.
\item Substitute $u=z+c$:
\[
L=\int_{-c}^0(z+c)^n\,dz=\frac{c^{n+1}}{n+1},\qquad
G=\int_0^{1-c}[1-(z+c)^n]\,dz
=\frac n{n+1}-c+\frac{c^{n+1}}{n+1}.
\]
Hence $\E[Z]=n/(n+1)-c$; fairness requires $c=n/(n+1)$.
$L$ is the expected size of losses, counting zero on profitable outcomes;
$G$ is the expected size of gains, counting zero on losing outcomes.
Neither is a conditional average given a loss or gain.
\item $\P(Z<0)=c^n$ and
$\E[|Z|]=G+L=n/(n+1)-c+2c^{n+1}/(n+1)$.
The absolute value before averaging removes cancellations between losses
and gains. At a fair price, the signed mean is zero but the absolute mean
is positive.
\item Mean profit: $0$. Loss probability: $(3/4)^3=27/64$.
Mean absolute payoff: $2(3/4)^4/4=81/512$ dollars.
Fairness is zero expected net payoff; it need not mean a one-half chance
of losing or an absence of risk.
\end{enumerate}
\end{solution}

\problem{A clock, a bus, and a normal measurement}{14}
\begin{enumerate}
\item A service has independent Poisson counts on disjoint intervals,
with mean $3t$ requests in $t$ minutes. Let $W$ be the time to the next
request. Derive its survival function from a zero-count probability;
find its mean, variance, and median. \partpts{4}
\item Compute $\P(W>5/3\mid W>1)$ and compare it with $\P(W>2/3)$.
A bus wait $V$ is instead uniform on $[0,2]$ minutes. Compute
$\P(V>5/3\mid V>1)$ and $\P(V>2/3)$; compare. \partpts{4}
\item The duration $X$ of a task is modeled as normal with mean 10 and
variance 9 minutes squared. Express $\P(8<X<13)$ and $\P(X>13)$ in
terms of the standard normal CDF $\Phi$. \partpts{3}
\item What density-height or endpoint mistake appears in
``$\P(X=10)=1/(3\sqrt{2\pi})$''? For $Y=2X-5$, give the distribution,
mean, and variance. \partpts{3}
\end{enumerate}
\begin{solution}
\begin{enumerate}
\item $\P(W>t)=\P(\Pois(3t)=0)=e^{-3t}$ for $t\ge0$.
Thus $W\sim\operatorname{Exp}(3)$, with mean $1/3$ minute,
variance $1/9$ minutes squared, and median $(\log2)/3$ minute.
For example, its mean is the tail integral $\int_0^\infty e^{-3t}\,dt$.
\item For the exponential, $e^{-5}/e^{-3}=e^{-2}=\P(W>2/3)$.
For the uniform, $(1/6)/(1/2)=1/3$, whereas
$\P(V>2/3)=2/3$. Only the exponential is memoryless.
\item The standard deviation is 3. The probabilities are
$\Phi(1)-\Phi(-2/3)$ and $1-\Phi(1)$.
\item A density height is not a point probability: $\P(X=10)=0$.
$Y\sim\operatorname{Normal}(15,36)$, where the second parameter is
variance. Its mean is $2(10)-5=15$ and variance is $2^2(9)=36$.
\end{enumerate}
\end{solution}

\par\addvspace{12pt}\Needspace{8\baselineskip}
{\bfseries\color{mint}Optional challenge: the rare sticker}\pts{10 extra}\par\nobreak
Now type $i$ has probability $p_i>0$, with $\sum_{i=1}^np_i=1$.
Purchases are still independent. Let $T$ count draws until all types appear.
Write $p_A=\sum_{i\in A}p_i$.
\begin{enumerate}
\item Explain why the next-new-type probability after $k$ types have
appeared can depend on \emph{which} types they were. Derive
$\P(T>m)=\sum_{\varnothing\ne A\subseteq\{1,\ldots,n\}}
(-1)^{|A|+1}(1-p_A)^m$ for $m\ge1$. \partpts{3}
\item Use the discrete tail sum to obtain
$\E[T]=\sum_{\varnothing\ne A}(-1)^{|A|+1}/p_A$.
For three types with probabilities $1/2,1/3,1/6$, evaluate it and
compare with three equally likely types. \partpts{4}
\item For two types with probabilities $p,1-p$, show the mean is
$1/p+1/(1-p)-1$ and is minimized uniquely at $p=1/2$. \partpts{3}
\end{enumerate}
\begin{solution}
\begin{enumerate}
\item After seeing the set $S$, the new-type probability is
$1-\sum_{i\in S}p_i$, which need not be determined by $|S|$.
A specified nonempty set $A$ is missing after $m$ draws with probability
$(1-p_A)^m$; inclusion--exclusion for the union of these missing events
gives the tail. Uniform geometric-stage parameters no longer apply.
\item The $m=0$ tail is 1. The alternating subset coefficients also sum
to 1, so we can include that term with the convention that each set is
missing before any draw (including the full set).
There are finitely many subsets and every $p_A>0$, so their convergent
geometric series may be summed separately:
$\sum_{m\ge0}(1-p_A)^m=1/p_A$.
For three types,
\[
\E[T]=2+3+6-\left(\frac65+\frac32+2\right)+1
=\frac{73}{10}=7.3,
\]
compared with $3H_3=11/2=5.5$ for uniform types.
\item $\E[T]=1/[p(1-p)]-1\ge4-1=3$, since
$p(1-p)=1/4-(p-1/2)^2\le1/4$. Equality holds exactly at $p=1/2$.
\end{enumerate}
\end{solution}
\end{document}
